\Chapter{Összegzés}

A dolgozat célja a Steam platformról származó videójáték-adatok
elemzése és különböző osztályozási feladatok megfogalmazása volt
gépi tanulási módszerek alkalmazásával.
A kutatás során egy valós, heterogén szerkezetű adathalmaz került
feldolgozásra, amely numerikus változókat,
kategória jellegű metaadatokat és szöveges leírásokat is tartalmazott.
A munka egyik fontos célkitűzése az volt,
hogy bemutassa, miként használhatók fel ezek az eltérő típusú információk
prediktív modellezési feladatokban.

A feldolgozás első lépéseként az adatok tisztítása és egységesítése történt meg.
Ez magában foglalta a hiányzó értékek kezelését,
a változók típusának egységesítését,
valamint a modellek számára megfelelő jellemzők kialakítását.
A kategória jellegű mezők bináris reprezentációvá alakítása,
illetve a szöveges leírások TF--IDF alapú vektorizálása
lehetővé tette, hogy az eltérő típusú adatok közös jellemzőtérben jelenjenek meg.

A dolgozat két eltérő osztályozási problémát vizsgált.
Az első feladat a videójátékok sikerességének becslése volt
a felhasználói értékelésekből származtatott címke alapján.
Ebben az esetben a strukturált metaadatok és numerikus változók
alkották a prediktív modellek bemenetét.
Az elemzés rámutatott arra, hogy a játékok egyes funkcionális
jellemzői és platformszolgáltatásai kapcsolatban állhatnak
a felhasználói visszajelzések által tükrözött sikerességgel.

A második vizsgált feladat a multiplayer és singleplayer
játékok automatikus elkülönítése volt.
Ebben az esetben a szöveges leírások és a strukturált tagek
bizonyultak a leginformatívabb jellemzőknek.
A szöveges reprezentáció lehetővé tette,
hogy a modellek a játékmechanikákra és funkciókra utaló
kifejezések alapján különítsék el a két kategóriát.

A különböző osztályozási modellek összehasonlítása azt mutatta,
hogy mind a lineáris, mind a nemlineáris módszerek
képesek értelmezhető és stabil predikciókat adni.
A lineáris modellek előnye elsősorban az értelmezhetőségben
és a nagy dimenziós jellemzőterek hatékony kezelésében jelentkezett,
míg az ensemble módszerek bizonyos esetekben
jobb prediktív teljesítményt mutattak.
Az eredmények összességében azt jelzik,
hogy a videójátékokhoz kapcsolódó metaadatok
és szöveges információk alkalmasak
gépi tanulási modellek bemeneteként történő felhasználásra.

A dolgozat egyik fontos tanulsága,
hogy a strukturált és a szöveges jellemzők kombinációja
jelentősen növelheti az osztályozási modellek információtartalmát.
A videójátékok leírásai és metaadatai nemcsak marketingcélú
információkat hordoznak,
hanem a játékmechanikákra és felhasználói élményre
vonatkozó implicit jelzéseket is tartalmaznak,
amelyek a prediktív modellek számára is hasznosak lehetnek.

A bemutatott megközelítés több irányban is továbbfejleszthető.
A jövőbeni kutatások egyik lehetséges iránya
összetettebb természetesnyelv-feldolgozási módszerek alkalmazása lehet,
például szóbeágyazások vagy neurális nyelvi modellek bevonásával.
Ezek képesek lehetnek a szövegek mélyebb szemantikai szerkezetének
megragadására, ami tovább javíthatja az osztályozási teljesítményt.

További fejlesztési lehetőséget jelenthet
összetettebb gépi tanulási modellek alkalmazása is,
például mély neurális hálózatok vagy fejlettebb ensemble módszerek használata.
Emellett érdemes lehet a sikeresség fogalmának finomabb modellezése,
például többosztályos vagy regressziós megközelítés formájában,
ahol a célváltozó nem bináris,
hanem folytonos vagy több kategóriából áll.

Összességében a dolgozat bemutatta,
hogy a Steam platformon elérhető videójáték-adatok
alkalmasak adatvezérelt elemzések és prediktív modellezési feladatok
megvalósítására.
Az alkalmazott módszerek nemcsak a prediktív teljesítmény szempontjából
bizonyultak hasznosnak,
hanem lehetőséget adtak a videójátékok tulajdonságai
és a felhasználói visszajelzések közötti összefüggések
jobb megértésére is.
